% Source: Independent quadratic expansion; user hull model and KACTL.
\subsection{斜率优化：把转移化成直线}
\label{dp:slope-basics}
\subsubsection*{原理}

先把转移展开为
\[
 f_i=A_i+\min_{j<i}(k_jx_i+b_j),
\]
其中 $A_i$ 只依赖 i，每个旧状态 j 对应一条直线。这里 k、b 只由已经求好的旧状态 j 决定，x 只由当前 i 决定；若系数还同时依赖 i、j 而不能拆开，就不能直接套这个模型。查询横坐标 $x_i$ 的最小值，得到 $f_i$ 后再插入由 i 生成的新直线，保证只使用 $j<i$。凸包维护的是候选直线在每个 x 处的最优值。本文把直线的最大值包络称为上凸壳、最小值包络称为下凸壳；只需插入直线与查询最值，不必输出几何顶点。

两份模板的原生查询均为最大值。最小值使用恒等式
\[
 \min_j(k_jx+b_j)=-\max_j((-k_j)x-b_j).
\]
也就是\textbf{插入时斜率、截距同时取负，查询结果再取负，x 不变}。不要只取负截距，也不要把 query 的结果直接当最小值。取负后还需检查单调模板的斜率条件。

例如 $S_i=\sum_{t=1}^i a_t$，考虑
\[
 f_i=\min_{j<i}\{f_j+(S_i-S_j)^2+C\}.
\]
展开后为
\[
 f_i=S_i^2+C-\max_{j<i}\{2S_jS_i-(f_j+S_j^2)\}.
\]
因此插入斜率 $2S_j$、截距 $-(f_j+S_j^2)$，在 $x=S_i$ 查询上凸壳。初值 $f_0=S_0=0$，先插入 (0,0)；第 i 轮先 query 求出 $f_i$，再 add\_line 加入 i 的直线。这个例子已经把 min 改写成 max，\textbf{不要再次取负转换}。

若 $a_i$ 非负，斜率和查询横坐标均不减，可用单调凸包达到 $O(n)$。顺序任意时，用动态直线凸包；若横坐标值域预先固定，亦可复用第~\ref{ds:li-chao}~节李超树。若只保证斜率有序而查询无序，不能弹掉队首，应在完整凸包上二分。

凸包中的“斜率优化”并不要求浮点斜率。交点比较可以交叉相乘，整数查询的分界点可以向下取整；前者要控制乘积范围，后者要正确处理负数除法与同斜率。
